“Laat los wat je griefde.
Want de zusterliefde

is onvoorwaardelijk
en onophoudelijk.

Ik ben jouw bloedverwant!”
Emmy’s echte vijand

bleek dichterbij te zijn,
via de schone schijn,

dan ze ooit had gedacht.
Emma greep steeds de macht

door te beschuldigen
of vals te huldigen

als aanzienuitslover.
“Emmy leert veel over

wat volwassenen doen,
volgend op wat gezoen

en wat knuffelplezier,
via een kleine kier

in onze bedstede”,
luidde de lofrede.

Haar ma vond dat niet tof,
gaf slaag in plaats van lof

en negeerde het feit
dat Emma haar bedtijd

in de late uren
blijkbaar ook met gluren

onopgemerkt vulde.
Emma kreeg slechts hulde.

Emmy kreeg eerst een lach
en dat werd plots beklag.

De emotieslinger
werd door de indringer

abrupt aangezwengeld.
Macht binnengehengeld

met valse ontkrachting
door gevoelontkenning.

Want juist een erkenning
geeft je verademing

en ontneemt het zijn kracht.
Echter, Emma zocht macht

in manipulatie
als Machiavelli.

Doch de echte vijand
zat in haar binnenkant.

Chaos versus orde.
Loki nam de horde

over elk streng gebod
als sluwe Noorse God.

De held als wel schavuit
leefde zich stiekem uit.

Als schaduw verweesd
in Emmy’s vlakke geest.

Een nieuwe methode
met cijfers als code

had zijn weg gevonden.
De schooljuffen konden

zo overzicht houden.
De getallen zouden

leerlingen opdrijven
en domheid inwrijven.

Emma had weer een knul,
die dacht vanuit zijn lul.

Hij hielp zijn minnares
bij de Latijnse les

door stil voor te leggen
wat Emma moest zeggen.

De begijnse schooljuf
tartte hun beiden bluf

door glashard te stellen
dat hij bij het spellen

alles had voorgezegd.
Het werd door hem weerlegd

dat zij dat had gehoord.
Zijn woord tegen haar woord.

“Ga maar naar de abdis,
die gelooft je gewis!”

Hij biechtte meteen op.
“Je krijgt niet op je kop,

maar zelfs een goed cijfer,
voor de extra ijver

en de samenwerking.”,
was de jufs opmerking.

De biecht maakt immers schuld
tot iets wat wordt geduld.

“Het was gul bovendien.
Ik geef jullie een tien!”

Emma en haar minnaar
keken verrast naar haar

en toen blij naar elkaar.
“Kijk verder zelf maar,

hoe je de tien verdeelt!”
Het tweetal keek bespeeld,

maar vond het verdict fair.
De juf plaagde met flair:

“Hoe je de 1 en 0,
zowel eerlijk als gul,

onder elkaar verdeeld.”
Ze waren uitgespeeld.

Emma wou controle.
Zo werd haar frivole

geneigdheid verzwegen
en kon ze die plegen

door je af te leiden.
Ze kon je misleiden,

maar als je dat doorzag
dan raakte ze van slag

en kaatste ze de bal.
Haalde ze iets van stal

wat ze over je wist
als een afleidingslist

om je te verblinden.
“Zoekt en u zult vinden.”

Met jouw schuld vastgesteld
had zij het wisselgeld

om hetzelfde te doen.
Weg was al haar ‘fatsoen’.

Ze verborg gevoelens
van passie voor jongens

zoals moeders dat deed.
Vooraf een plan gesmeed

was moeder voorbereid
op elk logisch verwijt

door een schuldverschuiving
naar de ellendeling,

die haar ontmaskerde
en zo vernederde.

Vaak haalde Emma rap
de krenten uit de pap

om deze minder heet
op te dienen als leed.

Soms had vader het door
en gaf hij geen gehoor

aan manipulatie
ten koste van Emmy

door beschuldigfraude.
“We halen geen oude

koeien meer uit de sloot”,
was wat vader gebood.

Pa was niet verstoken
om zelf te gaan stoken

als wraak op de arglist.
Hij zaaide vaak een twist

als weer twee roddelaars
kropen in elkaar aars.

“Maar zonet zei je nog,
‘Tante is vol bedrog’!”

De oorlog ontketend
keek hij van niets wetend.

Emma vormde bondjes.
Haar vele schoothondjes

schermden haar ondeugd af.
De kudde was te laf

om elkaar openlijk,
constructief en eerlijk

op iets aan te spreken.
Zo leken haar streken

verborgen te blijven.
Want de oude wijven

wisten maar al te goed
‘wie zoet doet, zoet ontmoet’.

Wellicht kon men niets doen
in plaats van uit ‘fatsoen’

achter iemand zijn rug
te steken als een mug.

Zo konden de schapen
rustig blijven slapen.

Alsof men niets fout deed,
wat men niet dromend weet.

Alsof een slachtoffer
steeds in hun reiskoffer

door het zware leven
spits is meegegeven.

De dader in aantocht
moest nog worden gezocht.

“Ik ga je vermoorden!”
Echter deze woorden

konden het niet klaren.
Emmy trok aan de haren

uit haar pure wanhoop.
Emma had dit verloop

al volgaarne voorzien,
maar huilde wel nadien.

Zo groeide de schande,
die op Emmy landde.

“Zussen vormen één paar
en houden van elkaar.

Leer van de parabel
over Kaïn en Abel.

Wees dus niet afgunstig.
Dat stemt God niet gunstig.

Ook al houden we meer
van Emma’s zedenleer,

maak geen strijd om de kroon.
Want daar is het gewoon

dat vaak broers en zussen
elkaars vuur uitblussen”,

blaatte uit moeders mond
toen ze Emmy vastbond

als een geit aan een paal.
Normaal voor een etmaal,

maar het was zaterdag,
al wat laat op de dag.

Morgen was er de kerk.
Al vroeg het zondagswerk.

Het touw was lang genoeg
dat als de drang toesloeg

ze net bij de privaat
kon voor een grote daad.

“Als men als een bruut klaagt
en zich als beest gedraagt,

dan moet men maar leven
als hun is gegeven.”

zei moeder als uitvlucht
voor haar platte wraakzucht.

Emma uitte haar wrok.
‘Wijs naar een zondebok’

als haar mechanisme.
Girard’s mimetische

rivaliteit in beeld.
Emmy was uitgespeeld.

Vloeide Emma's misdaad
voort uit angst voor verraad?

Of steeg die bezorgdheid
door haar eigen valsheid

onbewust te leren
en te projecteren?
